\label{sec:background}

\subsection{The Historical Development of Incumbency as a Redistricting Criterion}

The use of incumbent home addresses in redistricting predates modern civil rights law.
Reconstruction-era state legislatures routinely drew districts to protect sitting members,
and this practice was so thoroughly embedded in the redistricting tradition that early
federal reapportionment litigation rarely addressed it~\citep{backstrom1978problems}.
The Supreme Court's entry into redistricting in \textit{Baker v.\ Carr}\footnote{
\textit{Baker v.\ Carr}, 369 U.S.\ 186 (1962).} and \textit{Reynolds v.\ Sims}
\footnote{\textit{Reynolds v.\ Sims}, 377 U.S.\ 533 (1964).} established the one-person,
one-vote principle but said nothing about whether mapmakers could consider incumbent
addresses in drawing equal-population districts.

The Court's first significant engagement with incumbency as a redistricting criterion
came in \textit{Karcher v.\ Daggett} (1983), where New Jersey's partisan gerrymander
was challenged on equal population grounds. Justice Brennan's majority opinion, striking
down the New Jersey map, noted that minor population deviations might be justified by
``legitimate objectives,'' including ``making districts compact, respecting municipal
boundaries, preserving the cores of prior districts, and avoiding contests between
incumbents.'' This list was explicitly non-exhaustive; the Court was cataloguing
justifications that had historically been accepted, not establishing a definitive taxonomy
of permissible criteria.

\subsection{Permissive vs.\ Mandatory: The Doctrinal Distinction}

The crucial doctrinal distinction that the incumbency objection typically elides is
between \textit{permissible} criteria (which a mapmaker may consider) and
\textit{mandatory} criteria (which a mapmaker must consider). The Supreme Court has
been consistent in characterising incumbency protection as permissible. It has never
held that any redistricting body is constitutionally required to consider incumbent
home addresses or to draw districts that place each incumbent in a protected home seat.

This distinction was clarified in \textit{Cox v.\ Larios}\footnote{\textit{Cox v.\
Larios}, 542 U.S.\ 947 (2004), summarily affirming 300 F.Supp.2d 1320 (N.D.\ Ga.\
2004).} where the Court affirmed that incumbency protection, while legitimate, cannot
justify population deviations that would otherwise fail constitutional scrutiny.
The Court's reasoning implicitly confirms the permissive character of the criterion:
a mandatory criterion would be a reason to accept whatever population deviation it
required, but a permissible criterion is only a potential justification for deviations
that are otherwise within constitutional tolerance.

\subsection{\textit{Burdick v.\ Takushi} and Election Administration Stability}

\textit{Burdick v.\ Takushi}\footnote{\textit{Burdick v.\ Takushi}, 504 U.S.\ 428
(1992).} addressed Hawaii's prohibition on write-in voting. Justice White's majority
opinion established the sliding-scale balancing test for election law challenges:
courts weigh the state's asserted interest against the burden on First and Fourteenth
Amendment rights, with more severe restrictions requiring stronger justifications.
Justice White identified ``election administration'' stability as a cognisable state
interest, and subsequent lower courts have cited this language to support the
proposition that incumbent protection---which promotes continuity in representation
and reduces the disruption of primary contests between established officeholders---
serves a state interest in election administration stability~\citep{pildes2002foreword}.

The application of \textit{Burdick} to redistricting is strained. \textit{Burdick}
addressed a ballot restriction affecting individual voters, not a legislative line-drawing
choice. But the lower court adoption of \textit{Burdick}'s interest-balancing framework
for redistricting incumbency questions means that practitioners must address it.
The response is straightforward: even accepting the \textit{Burdick} framework,
the state's interest in election administration stability does not require that every
redistricting plan protect every incumbent. The interest is served by having some
continuity in representation, not by guaranteeing every sitting member a safe home
district.

\subsection{Summary of Federal Doctrine}

Table~\ref{tab:federal-cases} summarises the federal case law on incumbency as a
redistricting criterion.

\begin{table}[ht]
\centering
\begin{threeparttable}
\caption{Federal case law on incumbency as a redistricting criterion.}
\label{tab:federal-cases}
\begin{tabular}{p{4.5cm} l p{6cm}}
\toprule
Case & Year & Holding on Incumbency \\
\midrule
\textit{Karcher v.\ Daggett} & 1983 & Avoiding incumbent contests is a ``legitimate
  objective'' that may justify minor pop.\ deviations; not mandatory \\
\addlinespace
\textit{Burdick v.\ Takushi} & 1992 & Election administration stability is a cognisable
  state interest; indirect support for incumbency consideration \\
\addlinespace
\textit{Cox v.\ Larios} & 2004 & Incumbency protection cannot justify pop.\ deviations
  exceeding constitutional tolerance; confirms permissive (not mandatory) character \\
\addlinespace
\textit{Rucho v.\ Common Cause} & 2019 & Partisan gerrymandering non-justiciable
  federally; state courts may evaluate incumbency as tool of partisan advantage \\
\bottomrule
\end{tabular}
\begin{tablenotes}
\small
\item Federal doctrine is consistently permissive: incumbency may be considered, not
  that it must be. No federal court has required a redistricting plan to protect
  any specific incumbent.
\end{tablenotes}
\end{threeparttable}
\end{table}
